% Vietnamese translation for OI Public Library
% Source-PDF: IOI中国国家候选队论文/2009/周而进/浅谈估价函数在信息学竞赛中的应用.pdf
\documentclass[11pt]{article}
\usepackage{oipl-vn}

\newcommand{\Oh}{\mathcal{O}}
\newcommand{\dist}{\operatorname{dist}}
\newcommand{\opt}{\operatorname{opt}}
\newcommand{\ANS}{\operatorname{ANS}}
\newcommand{\MST}{\operatorname{MST}}

\title{Bàn sơ về ứng dụng hàm đánh giá trong thi tin học}
\author{Zhou Erjin\\Trường Trung học số 1 Thiệu Hưng, Chiết Giang}
\date{2009}

\begin{document}
\maketitle

\origtitle{Bàn sơ về ứng dụng hàm đánh giá trong thi tin học}
\sourcepath{IOI China candidate team papers / 2009 / Zhou Erjin}

\begin{abstract}
Hàm đánh giá có vai trò quan trọng trong một lớp bài toán tinh chỉnh dần lời
giải, nơi không gian trạng thái rất lớn và chuyển trạng thái tốn thời gian.
Bài viết bàn về ứng dụng của hàm đánh giá trong thi tin học từ ba góc độ:
tìm kiếm, tối ưu quy hoạch động và tìm lời giải khá tốt. Qua đó trình bày
cách chọn, đánh giá và sử dụng hàm đánh giá để xử lý những bài vốn khó bắt
đầu.
\end{abstract}

\paragraph{Từ khóa.}
Hàm đánh giá; tinh chỉnh dần; tìm kiếm; quy hoạch động; lời giải khá tốt; cân
bằng.

\section{Mở đầu}

Hàm đánh giá là một cách ánh xạ mức tốt xấu của trạng thái trong tương lai về
một giá trị cụ thể. Ví dụ đơn giản là trong cờ vua: người mới chơi thường được
cho biết tốt \(1\) điểm, tượng hoặc mã \(3\) điểm, xe \(5\) điểm, hậu \(9\)
điểm. Với tiêu chuẩn đó, có thể xây dựng một đánh giá rất đơn giản: lấy giá
trị quân có thể ăn ở nước kế tiếp làm giá trị đánh giá cho nước đó. Nếu đưa
hàm đánh giá như vậy vào chương trình, chương trình vốn chỉ biết vét cạn đã
có phần nào năng lực của người chơi nhập môn.

Đề thi tin học thường biến hóa đa dạng. Khi chưa tìm được phương pháp tối ưu,
việc dùng hàm đánh giá hợp lý đôi khi đem lại hiệu quả rất bất ngờ.

\section{Hàm đánh giá trong tìm kiếm}

\subsection{Bài toán tối ưu}

Trong thi tin học ta thường gặp các bài tối ưu mà bản thân bài toán khó có
lời giải đa thức hiệu quả, nên phải dùng tìm kiếm. Không gian trạng thái của
loại bài này thường rất lớn; vét cạn trực tiếp chắc chắn kém hiệu quả. Nếu sử
dụng hàm đánh giá, ta có thể cắt nhánh đúng lúc và nâng cao hiệu suất chương
trình.

\subsubsection{Mines For Diamonds}

\paragraph{Mô tả bài toán.}
Cho một mỏ kim cương kích thước \(n\times m\). Cần đào một số đường hầm để lấy
hết kim cương. Một đường hầm là một đường gấp khúc xuất phát từ biên; các viên
kim cương nằm trên đường đều được lấy. Hai đường gấp khúc không được cắt nhau,
nhưng có thể kề nhau. Hãy tìm tổng độ dài đường hầm nhỏ nhất. Giới hạn
\(n,m\le 11\), số kim cương không quá \(10\).

\paragraph{Phân tích.}
Vì phạm vi nhỏ và không thấy phương pháp đa thức rõ ràng, ta nghĩ tới tìm
kiếm. Vét cạn trực tiếp sẽ TLE, nên cần cắt nhánh.

Khi độ sâu tăng, số nút sinh ra tăng cực nhanh. Nếu ở giai đoạn sớm có thể
cắt bỏ một phần trạng thái thì hiệu quả sẽ tăng mạnh. Cách tự nhiên là dùng
hàm đánh giá để cắt nhánh: nếu \(g+h\ge \mathit{ans}\), không cần tiếp tục tìm
kiếm. Ở đây \(g\) là chi phí đã dùng để tới trạng thái hiện tại, \(h\) là ước
lượng lạc quan từ trạng thái hiện tại tới mục tiêu. Điều kiện cần là
\[
  h\le h^*,
\]
trong đó \(h^*\) là chi phí thật sự còn cần.

Vấn đề chính là thiết kế \(h\). Quan sát đề, ràng buộc khó là hai đường hầm
không được cắt nhau. Nếu bỏ qua ràng buộc này, bài toán chuyển thành: cho một
tập điểm, chia chúng thành vài tập con; trong mỗi tập con nối các điểm bằng
một dây chuyền sao cho tổng chi phí nhỏ nhất. Bài này có thể giải bằng quy
hoạch động hai tầng. Gọi \(\opt[S]\) là lời giải tối ưu cho tập điểm \(S\), ta
có:
\[
  \opt[S]=\min_{T\subset S}\{\opt[S-T]+w[T]\}.
\]
Ở đây \(w[T]\) là chi phí nhỏ nhất để nối các điểm trong \(T\) bằng một dây
chuyền:
\[
  w[T]=\min_{k\in T}\{w[T-\{k\}]+\dist[T-\{k\}][k]\}.
\]
Lấy \(\opt[S]\) làm hàm đánh giá vừa thỏa \(h\le h^*\), vừa thường rất gần
chi phí tối ưu thực tế. Trong dữ liệu thử của bài gốc, số nút sinh ra giảm từ
mức hàng chục triệu xuống rất nhỏ, và chương trình có thể AC trên UVa trong
dưới \(0.1\) giây.

\subsubsection{Jaguar King}

\paragraph{Mô tả bài toán.}
Có \(n\) con báo xếp hàng ở các vị trí \(1..n\), \(n\) chia hết cho \(4\). Báo
số \(1\) là vua báo; chỉ vua báo được đổi chỗ với báo khác. Nếu vua đang ở vị
trí \(i\), các nước nhảy hợp lệ phụ thuộc vào \(i\bmod 4\):
\[
\begin{array}{c|l}
i\bmod4 & \text{vị trí có thể nhảy tới}\\
\hline
1 & i+1,\ i+3,\ i-4,\ i+4\\
2 & i+1,\ i-1,\ i-4,\ i+4\\
3 & i+1,\ i-1,\ i-4,\ i+4\\
0 & i-3,\ i-1,\ i-4,\ i+4
\end{array}
\]
Hỏi vua báo cần nhảy ít nhất bao nhiêu lần để xếp dãy đúng thứ tự. Giới hạn
\(n\le40\).

\paragraph{Phân tích.}
Bài toán giống mô hình \(8\)-puzzle: chỉ một vị trí có thể di chuyển. Nếu sắp
xếp các vị trí theo dạng bảng có \(4\) cột, ta thấy mỗi điểm thực chất có thể
nhảy tới bốn ô kề. Vì thế có thể dùng đánh giá tương tự khoảng cách Manhattan:
\[
  E(s)=\sum_{i>1}\dist(i,\mathit{pos}[i]),
\]
trong đó \(\dist(i,j)\) là khoảng cách Manhattan trên bảng \(4\) cột:
\[
  x=\left|\left\lfloor\frac{i-1}{4}\right\rfloor
      -\left\lfloor\frac{j-1}{4}\right\rfloor\right|,\qquad
  y=|(i-1)\bmod4-(j-1)\bmod4|,
\]
nếu \(y=3\) thì đặt \(y=1\), và \(\dist(i,j)=x+y\). Thêm đánh giá này vào tìm
kiếm sâu lặp sẽ vượt qua dữ liệu UVa nhanh chóng.

\paragraph{Mở rộng.}
Nếu mỗi điểm \(V_i\) có một tập điểm tới được \((V_{i1},\ldots,V_{ik})\) với
chi phí \((W_{i1},\ldots,W_{ik})\), không còn quan hệ Manhattan. Khi đó có thể
dựng đồ thị \(G=(V,E)\), với cạnh tương ứng và trọng số tương ứng, rồi dùng
khoảng cách đường ngắn nhất:
\[
  E(s)=\sum_{i>1}\mathit{ssp}[i][\mathit{pos}[i]].
\]
Ở đây \(\mathit{ssp}[i][j]\) là đường ngắn nhất từ \(i\) tới \(j\) trong đồ thị
\(G\).

\subsubsection{Tiểu kết}

Với bài tối ưu dùng tìm kiếm, tìm kiếm mù thường duyệt quá nhiều trường hợp.
Đánh giá kịp thời trạng thái trở thành yếu tố quan trọng. Hàm đánh giá hoàn
hảo sẽ bằng đúng chi phí thật còn lại, nhưng khó có ``lời tiên tri'' như vậy.
Trong thực tế, ta thường bỏ qua một ràng buộc nào đó của đề, biến bài toán
thành một bài dễ hơn, rồi dùng lời giải của bài dễ hơn làm cận. Bài dễ hơn
thường giải được bằng tham lam hoặc quy hoạch động, từ đó cắt bỏ rất nhiều
trạng thái vô ích.

\subsection{Bài toán khả thi}

Có một lớp bài tìm kiếm không yêu cầu lời giải tối ưu, chỉ cần tìm một lời giải
hợp lệ. Không gian trạng thái của chúng thường còn lớn hơn; cắt nhánh tối ưu
kiểu thường dùng có thể làm hiệu quả giảm mạnh. Một cách hiệu quả là dùng hàm
đánh giá để phân biệt trạng thái nào đáng thử trước.

\subsubsection{Sokoban}

\paragraph{Mô tả bài toán.}
Cho một thế Sokoban kích thước \(n\times m\), cần tìm một lời giải không quá
\(10000\) bước. Dữ liệu bảo đảm có lời giải, \(n,m\le8\). Ký hiệu:
\begin{itemize}
  \item \texttt{\#}: tường;
  \item \texttt{.}: đích;
  \item \texttt{@}: vị trí ban đầu của người;
  \item \texttt{+}: người đứng trên đích;
  \item \texttt{\$}: hộp;
  \item \texttt{*}: hộp đứng trên đích.
\end{itemize}

\paragraph{Phân tích.}
Diện tích thực sự tối đa sau khi bỏ viền tường chỉ khoảng \(6\times6\), nên số
hộp tối đa khoảng \(35\). Vì có nhiều hộp, để dễ viết ta xem mỗi bước là đẩy
một hộp cho tới khi mọi hộp vào đích. Khó khăn lớn là deadlock: hộp bị đẩy vào
góc tường, bốn hộp dính nhau, hộp sát một cạnh tường không có đích, v.v. Nhiều
deadlock người nhìn ra ngay nhưng máy khó kiểm tra, còn kiểm tra thường xuyên
lại tốn thời gian.

Thay vì đẩy hộp, ta suy nghĩ ngược: biến bài toán thành kéo hộp. Như vậy các
deadlock nói trên tự nhiên biến mất.

Định nghĩa trạng thái \(S(P,Q,\mathit{pos})\), trong đó \(P\) là tập vị trí hộp,
\(Q\) là tập đích, và \(\mathit{pos}\) là vị trí người. Dùng tìm kiếm sâu lặp:
liệt kê giới hạn đáp án \(\mathit{MAX}\). Nếu chi phí hiện tại là \(g\), khi
\[
  g+h>\mathit{MAX}
\]
thì cắt nhánh.

Một đánh giá hiển nhiên là khoảng cách Manhattan từ mỗi hộp tới một đích gần
nhất:
\[
  h(S)=
  \sum_{i\in P}\min_{j\in Q}
  \bigl(|X_i-X_j|+|Y_i-Y_j|\bigr).
\]
Đánh giá này hợp lệ nhưng chưa đủ mạnh; vì nó ghép nhiều hộp vào cùng một đích
nên có thể đánh giá quá lạc quan. Cải tiến tự nhiên là dùng ghép cặp tối ưu
KM giữa hộp và đích. Dù KM có độ phức tạp \(\Oh(n^3)\), ở đây \(n\) nhỏ. Hơn
nữa, sau một lần đẩy chỉ có một hộp đổi vị trí, nên từ ghép cặp tối ưu cũ,
ghép cặp mới khác nhiều nhất hai cặp; có thể tối ưu cập nhật xuống gần
\(\Oh(n)\) sau lần đầu.

Hàm đánh giá không chỉ dùng để cắt nhánh. Nó còn so sánh độ hứa hẹn giữa các
con của một trạng thái. Nếu trước khi mở rộng một nút, ta sắp xếp các con theo
hàm đánh giá và tìm trạng thái hứa hẹn trước, chương trình sẽ tìm lời giải sớm
hơn nhiều. Trong bài gốc, sau tối ưu này mọi bộ thử đều chạy được và qua toàn
bộ dữ liệu Ural.

\paragraph{Lưu ý về tính đơn điệu.}
Hàm đánh giá trong Sokoban không đơn điệu, nghĩa là không bảo đảm lần đầu gặp
một trạng thái là đường ngắn nhất tới trạng thái đó. Với tìm kiếm sâu lặp và
bảng băm, điều này có thể làm mất một lời giải ở độ sâu hiện tại. Tuy nhiên
bài chỉ yêu cầu một lời giải khả thi; nếu ở tầng hiện tại bỏ lỡ, ở tầng sâu
hơn vẫn có thể tìm lại. Đổi lại, hàm đánh giá giúp hiệu quả tăng rất mạnh.

\subsubsection{Tiểu kết}

Với bài khả thi, hàm đánh giá phải thể hiện tốt sự khác biệt giữa các trạng
thái. Nó không chỉ cho biết trạng thái hiện tại có đáng tiếp tục hay không, mà
còn giúp so sánh hai trạng thái để nhanh chóng tìm được một lời giải. Tính
không đơn điệu của hàm đánh giá trong loại bài này thường có thể chấp nhận
được nếu đổi lại hiệu suất tốt hơn.

\section{Hàm đánh giá trong tối ưu quy hoạch động}

Quy hoạch động thường nổi tiếng vì gọn và hiệu quả. Nhưng nếu số trạng thái quá
lớn, quy hoạch động trực tiếp vẫn lực bất tòng tâm. Khi đó có thể dùng hàm
đánh giá để giảm số trạng thái hoặc ưu tiên thứ tự xử lý.

\subsection{Largest Fence}

\paragraph{Mô tả bài toán.}
Farmer John có \(n\) điểm hàng rào, với \(5\le n\le250\). Ông cần dựng một
vòng hàng rào là một bao lồi và muốn số điểm trên bao lồi là nhiều nhất.

\paragraph{Phân tích.}
Một lời giải \(\Oh(N^4)\) khá trực tiếp là: liệt kê điểm trái nhất trên bao
lồi cuối cùng, sắp mọi điểm bên phải theo góc cực, rồi dùng trạng thái
\(\opt[i][j]\) biểu thị ``bao lồi'' hiện tại có hai điểm cuối là \(i,j\). Chuyển
trạng thái:
\[
  \opt[i][j]=\max\{\opt[j][k]+1\},
\]
trong đó \(\operatorname{Point}(i)\) nằm bên trái đường thẳng
\(\operatorname{Line}(k,j)\).

Số trạng thái là \(\Oh(N^3)\), mỗi chuyển tốn \(\Oh(N)\), nên tổng là
\(\Oh(N^4)\). Cách này chưa đủ nhanh cho dữ liệu lớn.

\paragraph{Tối ưu 1.}
Nếu thiết kế được một hàm đánh giá \(h(s)\) để phán đoán trạng thái hiện tại có
thể dẫn tới nghiệm tốt hơn hay không, ta sẽ giảm được nhiều trạng thái. Quan
sát một bao lồi, từ một đỉnh trên bao lồi tới các đỉnh khác, dãy khoảng cách
thường có dạng đơn đỉnh: trước tăng rồi giảm. Điều này gợi ý bài toán dãy
không giảm rồi không tăng.

Gọi \(\mathit{down}[i]\) là độ dài dãy giảm dài nhất bắt đầu từ điểm \(i\), và
\(\mathit{up}[i]\) là độ dài dãy đơn đỉnh dài nhất bắt đầu từ \(i\):
\[
  \mathit{down}[i]=
  \max\{\mathit{down}[j]+1\mid j>i,\ \dist[j]<\dist[i]\},
\]
\[
  \mathit{up}[i]=
  \max\left(\mathit{down}[i],\
  \max\{\mathit{up}[j]+1\mid j>i,\ \dist[j]>\dist[i]\}\right).
\]
Đặt
\[
  h[i][j]=\mathit{up}[j].
\]
Nếu
\[
  \opt[i][j]+h[i][j]\le \ANS,
\]
thì trạng thái này không thể dẫn tới đáp án tốt hơn và không cần mở rộng.
Trong dữ liệu cực hạn \(N=250\), bước này giảm thời gian từ hơn \(10\) giây
xuống khoảng vài giây.

\paragraph{Tối ưu 2.}
Giống bài Sokoban, hàm đánh giá còn có thể so sánh độ tốt của các trạng thái.
Trong cùng một tầng chuyển, ưu tiên xử lý trạng thái con có đánh giá cao hơn.
Làm vậy giúp tìm nghiệm tốt sớm hơn; nghiệm hiện tại càng tốt thì cắt nhánh
càng mạnh. Sau tối ưu này, chương trình có thể giải dữ liệu cực hạn dưới một
giây và vượt qua dữ liệu USACO.

\subsection{Tiểu kết}

Quy hoạch động rất hiệu quả, nhưng khi số trạng thái khổng lồ và chuyển trạng
thái dày đặc, dùng trực tiếp thường khó đạt yêu cầu. Với các bài có dạng nhớ
hóa tìm kiếm hoặc trạng thái có thể đánh giá, hàm đánh giá giúp bỏ qua trạng
thái có triển vọng thấp hoặc ưu tiên trạng thái có triển vọng cao, nhờ đó nâng
cao hiệu suất đáng kể.

\section{Hàm đánh giá trong bài tìm lời giải khá tốt}

Một số bài trong thi đấu chỉ yêu cầu lời giải khá tốt, không cần tối ưu tuyệt
đối. Loại bài này linh hoạt hơn và kiểm tra năng lực phán đoán của thí sinh.
Chúng thường là bài NP-đầy đủ hoặc khó tìm lời giải đa thức trong thời gian thi.
Trong thực tế, dùng hàm đánh giá thường cho một thuật toán gọn và hiệu quả.

\subsection{Bài toán TSP}

\paragraph{Mô tả bài toán.}
Bài toán người bán hàng: cho \(n\) điểm trên mặt phẳng, tìm một chu trình đi
qua mỗi điểm đúng một lần và có tổng độ dài nhỏ nhất. Đây là bài NP-đầy đủ
kinh điển; với đồ thị tổng quát, việc tìm nghiệm tối ưu rất khó. Trong ứng dụng,
thường chỉ cần một nghiệm khá tốt.

\paragraph{Phân tích.}
Đối với một trạng thái \(s\), một đánh giá khả thi cần thỏa
\[
  h(s)\le h^*(s).
\]
Ta lại nới lỏng ràng buộc: đề yêu cầu một chu trình Hamilton, mỗi điểm có bậc
đúng bằng \(2\). Nếu bỏ qua ràng buộc đó, bài toán trở thành bài toán cây khung
nhỏ nhất. Cụ thể, gọi tập điểm là \(P\), và giả sử đã có một phần đường đi
\[
  Q=\{p_1,p_2,\ldots,p_k\}.
\]
Ta lấy tập điểm
\[
  P-Q+\{p_1,p_k\}
\]
và tính cây khung nhỏ nhất trên tập này. Chi phí của cây đó là cận dưới cho
phần đường còn lại.

\begin{center}
  \small Hình minh họa trong bản gốc: đường nét liền là phần đường đã chọn,
  đường nét đứt là cây khung nhỏ nhất dùng làm đánh giá cho phần còn lại.
\end{center}

Trong thử nghiệm của tác giả trên các bộ dữ liệu \(50..100\) điểm, chương
trình dùng đánh giá MST cho nghiệm tốt hơn rõ rệt so với điều chỉnh ngẫu nhiên.
Điều chỉnh ngẫu nhiên dễ rơi vào tối ưu cục bộ; dù có thể thêm biến dị kiểu
thuật toán di truyền, hiệu quả vẫn không ổn định bằng cách dùng hàm đánh giá.

\subsection{Tiểu kết}

Với bài chỉ cần nghiệm khá tốt, đề thường là NP-đầy đủ hoặc khó nghĩ ra lời
giải đa thức trong thời gian thi. Nếu trạng thái của bài tương đối rõ, ta có
thể chuyển bài gốc thành một bài dễ hơn để lấy cận trên hoặc cận dưới. Hàm
đánh giá trong lớp bài này có thể linh hoạt hơn, và cần cân bằng giữa thời
gian chạy với chất lượng nghiệm.

\section*{Tài liệu tham khảo}

\begin{enumerate}
  \item Frank Takes, \textit{Sokoban: Reversed Solving}.
  \item George F. Luger, \textit{Artificial Intelligence: Structures and
  Strategies for Complex Problem Solving}.
  \item Anand Panangadan, \textit{Incremental Evaluation of Minimum Spanning
  Trees for the Traveling Salesman Problem}.
\end{enumerate}

\section*{Lời cảm ơn}

Tác giả cảm ơn các thầy cô đã bồi dưỡng và hướng dẫn; cảm ơn các bạn cùng tập
luyện; và cảm ơn những người bạn trên mạng đã cùng trao đổi, thảo luận.

\end{document}
